\section{Sequencing the Compound Exposure Portfolio}
\label{sec:investment-sequencing}

\subsection{The Compound Priority Metric}

Sequencing the compound exposure investment portfolio requires a priority metric that reflects the combined urgency of both B1 isolation and D1 environmental vulnerability. We use the B1$\times$D1 product score: the product of the two dimension scores for each corridor, which rewards corridors where both dimensions are simultaneously elevated. A corridor with B1 = 8.3 and D1 = 7.8 (Donner Pass) scores 64.7 on this metric; a corridor with B1 = 5.0 and D1 = 5.0 scores 25.0 --- even though both sum to 13.0 and 10.0 respectively.

The compound priority sequence for the 11 identified corridors, ranked by B1$\times$D1 product:

\begin{center}
\begin{tabular}{llrrrr}
\toprule
\textbf{Rank} & \textbf{Corridor} & \textbf{B1} & \textbf{D1} & \textbf{Product} & \textbf{Phase} \\
\midrule
1 & I-80 Donner Pass (CA)       & 8.3 & 7.8 & 64.7 & 1 \\
2 & I-10 Gulf Coast LA          & 6.8 & 8.4 & 57.1 & 2 \\
3 & I-90 Snoqualmie Pass (WA)   & 7.6 & 7.1 & 54.0 & 3 \\
4 & I-5 Siskiyou Pass (OR)      & 7.1 & 6.8 & 48.3 & 3 \\
5 & I-35 Oklahoma Tornado Corr. & 6.2 & 7.2 & 44.6 & 4 \\
6 & I-90 Homestake Pass (MT)    & 7.8 & 6.4 & 49.9 & 4 \\
7 & I-10 Gulf Coast TX (Houston)& 5.9 & 7.6 & 44.8 & 4 \\
\bottomrule
\end{tabular}
\end{center}

\subsection{Why the Compound Ordering Differs from Single-Dimension Orderings}

The compound ordering produces materially different investment sequences than either B1-only or D1-only sequencing.

\textbf{B1-only ordering} would prioritize: Donner (8.3), Homestake (7.8), Snoqualmie (7.6), Siskiyou (7.1). This sequence deprioritizes Gulf Coast I-10 (B1 = 6.8) despite its \$820M annual disruption cost, because I-20 provides a functional (if costly) alternate. A B1-only investment program would invest in Homestake Pass before Gulf Coast I-10 --- but Homestake closes 25 times per year with low annual freight cost, while Gulf Coast closes 8 times with high freight cost because each closure is longer and affects more freight.

\textbf{D1-only ordering} would prioritize: Gulf Coast LA (8.4), Gulf Coast TX (7.6), Donner (7.8), I-35 Oklahoma (7.2), Snoqualmie (7.1). This sequence would invest in Gulf Coast LA before Donner, despite Donner's \$1.3B annual disruption cost exceeding Gulf Coast's \$0.82B. The reversal occurs because Gulf Coast LA has the highest environmental exposure (D1 = 8.4) even though its B1 isolation is moderate.

The compound metric correctly identifies Donner as the first priority because it is both severely isolated (B1 = 8.3) and frequently closed (D1 = 7.8). A dollar invested at Donner prevents more national freight disruption than a dollar invested at any single-dimension priority, precisely because both the structural and environmental conditions are simultaneously at their worst.

\subsection{Phased Investment Program}

\textbf{Phase 1 (immediate, FY2026--FY2029): Donner Freight Tunnel.}
The binding constraint for Phase 1 is political and financial, not technical: the Donner tunnel at \$4.0B exceeds available IIJA resilience program annual obligations for a single project. Phase 1 requires a dedicated authorization, analogous to the Surface Transportation Reauthorization Act mechanisms used for major bridge replacement projects (e.g., the \$3.3B authorization for the Portal Bridge replacement on the Northeast Corridor). Engineering feasibility work and environmental review can begin immediately on existing FHWA planning funds.

\textbf{Phase 2 (FY2027--FY2030): Gulf Coast I-10 hardening and alternate designation.}
The \$2.9B Gulf Coast portfolio (Louisiana segment hardening + I-20 freight alternate designation) is divisible into annual appropriation increments and fits within a single IIJA formula program cycle. The D1 trajectory for Gulf Coast I-10 is the steepest in the corpus: delay of 4 years increases the 2050 compound exposure significantly as sea-level-rise projections shift from mid-range to high-range scenarios.

\textbf{Phase 3 (FY2030--FY2035): Snoqualmie snowshed and Siskiyou hardening.}
Both projects are below \$2.0B and addressable within standard NHPP formula apportionments with supplemental discretionary grants under RAISE or INFRA programs. Snoqualmie snowshed construction has a precedent in the existing I-90 avalanche control infrastructure; Siskiyou hardening involves tunnel ventilation upgrades, chain-control improvement, and alternate route (OR-99) load rating improvement.

\textbf{Phase 4 (FY2035+): Remaining compound corridors.}
I-35 Oklahoma, Homestake Pass, and the Gulf Coast Texas segment (Houston). The I-35 Oklahoma investment (US-77 upgrade, \$1.8B) should be evaluated against the compound priority metric at the time of Phase 4 initiation, as tornado corridor D1 scores may shift with updated NOAA convective weather projections.

\subsection{Total Portfolio Envelope}

The full compound exposure investment program (Phases 1--4, top 7 corridors) requires approximately \$12.4B in capital investment and is projected to yield \$5.9B in annual avoided disruption costs at full build-out. At a 7\% discount rate, the 30-year portfolio NPV is approximately \$62B against \$12.4B invested --- a 5.0:1 cost-benefit ratio. The portfolio payback period at undiscounted cash flows is approximately 2.1 years from Phase 1 completion.
